% GENERATED FILE. Edit canonical lessons, then run npm run generate:books.
% canonical-source-hash: fnv1a64:9d3684e4f4399141
% canonical-lessons: TE-C536-vanabhojanam, TE-C536-kalaksepam, TE-C536-totapani, TE-C536-allika, TE-C536-pekata

\chapter{Free Time}
\label{ch:te-free-time}
\hlchaptermodality{536}

\begin{chapteropening}
\textbf{By the end of this chapter:} \emph{I can talk about picnics, pastimes, gardening, knitting and cards.}

Complete the last lesson of chapter 536: I can talk about picnics, pastimes, gardening, knitting and cards.
\end{chapteropening}

\section[\texorpdfstring{\te{వనభోజనం} (vanabhōjanaṁ)}{వనభోజనం (vanabhōjanaṁ)}]{\te{వనభోజనం} (vanabhōjanaṁ) — a picnic; a meal outdoors}

\label{lesson:TE-C536-vanabhojanam}



\begin{quote}
Before the new one: say the Telugu for to call to account, then the Telugu for to deceive.
\end{quote}

\subsection*{You'll want to know: \te{వనభోజనం}}
\textbf{\te{వనభోజనం}} — \emph{vanabhōjanaṁ} — \textquotedblleft{}a picnic; a meal outdoors\textquotedblright{}.

\begin{quote}
\te{ఆదివారం} \te{మేము} \te{స్నేహితులతో} \te{కలిసి} \te{వనభోజనానికి} \te{వెళ్ళాం}.
\end{quote}

\begin{quote}
\emph{ādivāraṁ mēmu snēhitulatō kalisi vanabhōjanāniki veḷḷāṁ.}
\end{quote}

\begin{quote}
On Sunday we went on a picnic with friends.
\end{quote}

\te{వనం}, a grove, + \te{భోజనం}, a meal: a meal eaten out under the trees.

\begin{cousinweb}
\textbf{\te{వనభోజనం}} (\emph{vanabhōjanaṁ}) is a learned compound built from Sanskrit \emph{vana}, \textquotedblleft{}a forest, a grove\textquotedblright{} (CDIAL 11258), and \emph{bhojana}, \textquotedblleft{}feeding, a meal\textquotedblright{} (Monier-Williams).
\end{cousinweb}

\begin{grammarlens}[title={What you've built}]
One more word for describing what you lived and felt.
\end{grammarlens}

\subsection*{Your turn}
\begin{itemize}
\raggedright
  \item \emph{Say it:} \emph{vanabhōjanaṁ}
  \item \emph{Say it:} \emph{vanabhōjanaṁ}, once more
  \item \emph{Say it:} \emph{ādivāraṁ mēmu snēhitulatō kalisi vanabhōjanāniki veḷḷāṁ.} — On Sunday we went on a picnic with friends.
  \item \emph{Recall:} say the Telugu for to call to account, then the Telugu for to deceive, then say \emph{vanabhōjanaṁ} again
\end{itemize}

\begin{tcolorbox}[breakable,colback=teal!4,colframe=teal!35!black,title={Before you move on}]
What does \te{వనభోజనం} mean? (\textquotedblleft{}A picnic; a meal outdoors\textquotedblright{}.) Say it once more.
\end{tcolorbox}
\section[\texorpdfstring{\te{కాలక్షేపం} (kālakṣēpaṁ)}{కాలక్షేపం (kālakṣēpaṁ)}]{\te{కాలక్షేపం} (kālakṣēpaṁ) — a pastime}

\label{lesson:TE-C536-kalaksepam}



\begin{quote}
Before the new one: say the Telugu for to deceive, then the Telugu for a picnic.
\end{quote}

\subsection*{You'll want to know: \te{కాలక్షేపం}}
\textbf{\te{కాలక్షేపం}} — \emph{kālakṣēpaṁ} — \textquotedblleft{}a pastime\textquotedblright{}.

\begin{quote}
\te{తాతయ్యకు} \te{పుస్తకాలు} \te{చదవడమే} \te{కాలక్షేపం}.
\end{quote}

\begin{quote}
\emph{tātayyaku pustakālu cadavaḍamē kālakṣēpaṁ.}
\end{quote}

\begin{quote}
Reading books is grandpa's pastime.
\end{quote}

Literally throwing time: passing the time pleasantly.

\begin{cousinweb}
\textbf{\te{కాలక్షేపం}} (\emph{kālakṣēpaṁ}) is a learned compound built from Sanskrit \emph{kāla}, \textquotedblleft{}time\textquotedblright{} (CDIAL 3084), and \emph{kṣēpa}, \textquotedblleft{}sending, throwing\textquotedblright{} (CDIAL 3738): sending the time on its way.
\end{cousinweb}

\begin{grammarlens}[title={What you've built}]
One more word for describing what you lived and felt.
\end{grammarlens}

\subsection*{Your turn}
\begin{itemize}
\raggedright
  \item \emph{Say it:} \emph{kālakṣēpaṁ}
  \item \emph{Say it:} \emph{kālakṣēpaṁ}, once more
  \item \emph{Say it:} \emph{tātayyaku pustakālu cadavaḍamē kālakṣēpaṁ.} — Reading books is grandpa's pastime.
  \item \emph{Recall:} say the Telugu for to deceive, then the Telugu for a picnic, then say \emph{kālakṣēpaṁ} again
\end{itemize}

\begin{tcolorbox}[breakable,colback=teal!4,colframe=teal!35!black,title={Before you move on}]
What does \te{కాలక్షేపం} mean? (\textquotedblleft{}A pastime\textquotedblright{}.) Say it once more.
\end{tcolorbox}
\section[\texorpdfstring{\te{తోటపని} (tōṭapani)}{తోటపని (tōṭapani)}]{\te{తోటపని} (tōṭapani) — gardening}

\label{lesson:TE-C536-totapani}



\begin{quote}
Before the new one: say the Telugu for a picnic, then the Telugu for a pastime.
\end{quote}

\subsection*{You'll want to know: \te{తోటపని}}
\textbf{\te{తోటపని}} — \emph{tōṭapani} — \textquotedblleft{}gardening\textquotedblright{}.

\begin{quote}
\te{ఆదివారాల్లో} \te{నాన్న} \te{తోటపని} \te{చేస్తారు}.
\end{quote}

\begin{quote}
\emph{ādivārāllō nānna tōṭapani cēstāru.}
\end{quote}

\begin{quote}
On Sundays dad does the gardening.
\end{quote}

\te{తోట}, garden, + \te{పని}, work.

\begin{cousinweb}
\textbf{\te{తోటపని}} (\emph{tōṭapani}) is a native compound: \emph{tōṭa}, \textquotedblleft{}garden, grove\textquotedblright{} (DEDR 3549), and \emph{pani}, \textquotedblleft{}work\textquotedblright{} (DEDR 3884).
\end{cousinweb}

\begin{grammarlens}[title={What you've built}]
One more word for describing what you lived and felt.
\end{grammarlens}

\subsection*{Your turn}
\begin{itemize}
\raggedright
  \item \emph{Say it:} \emph{tōṭapani}
  \item \emph{Say it:} \emph{tōṭapani}, once more
  \item \emph{Say it:} \emph{ādivārāllō nānna tōṭapani cēstāru.} — On Sundays dad does the gardening.
  \item \emph{Recall:} say the Telugu for a picnic, then the Telugu for a pastime, then say \emph{tōṭapani} again
\end{itemize}

\begin{tcolorbox}[breakable,colback=teal!4,colframe=teal!35!black,title={Before you move on}]
What does \te{తోటపని} mean? (\textquotedblleft{}Gardening\textquotedblright{}.) Say it once more.
\end{tcolorbox}
\section[\texorpdfstring{\te{అల్లిక} (allika)}{అల్లిక (allika)}]{\te{అల్లిక} (allika) — knitting; weaving}

\label{lesson:TE-C536-allika}



\begin{quote}
Before the new one: say the Telugu for a pastime, then the Telugu for gardening.
\end{quote}

\subsection*{You'll want to know: \te{అల్లిక}}
\textbf{\te{అల్లిక}} — \emph{allika} — \textquotedblleft{}knitting; weaving\textquotedblright{}.

\begin{quote}
\te{అమ్మమ్మ} \te{ఖాళీ} \te{సమయంలో} \te{అల్లిక} \te{చేస్తుంది}.
\end{quote}

\begin{quote}
\emph{ammamma khāḷī samayaṁlō allika cēstundi.}
\end{quote}

\begin{quote}
Grandma knits in her free time.
\end{quote}

The noun of \te{అల్లు}, to knit.

\begin{cousinweb}
\textbf{\te{అల్లిక}} (\emph{allika}) is native Dravidian: DEDR 260 lists Telugu \emph{allu}, \textquotedblleft{}to plait, to braid, to weave\textquotedblright{}, with Tamil \emph{allu}, \textquotedblleft{}to knit\textquotedblright{}.
\end{cousinweb}

\begin{grammarlens}[title={What you've built}]
One more word for describing what you lived and felt.
\end{grammarlens}

\subsection*{Your turn}
\begin{itemize}
\raggedright
  \item \emph{Say it:} \emph{allika}
  \item \emph{Say it:} \emph{allika}, once more
  \item \emph{Say it:} \emph{ammamma khāḷī samayaṁlō allika cēstundi.} — Grandma knits in her free time.
  \item \emph{Recall:} say the Telugu for a pastime, then the Telugu for gardening, then say \emph{allika} again
\end{itemize}

\begin{tcolorbox}[breakable,colback=teal!4,colframe=teal!35!black,title={Before you move on}]
What does \te{అల్లిక} mean? (\textquotedblleft{}Knitting; weaving\textquotedblright{}.) Say it once more.
\end{tcolorbox}
\section[\texorpdfstring{\te{పేకాట} (pēkāṭa)}{పేకాట (pēkāṭa)}]{\te{పేకాట} (pēkāṭa) — a card game}

\label{lesson:TE-C536-pekata}



\begin{quote}
Before the new one: say the Telugu for gardening, then the Telugu for knitting.
\end{quote}

\subsection*{You'll want to know: \te{పేకాట}}
\textbf{\te{పేకాట}} — \emph{pēkāṭa} — \textquotedblleft{}a card game\textquotedblright{}.

\begin{quote}
\te{పండుగ} \te{రాత్రి} \te{పెద్దలు} \te{పేకాట} \te{ఆడారు}.
\end{quote}

\begin{quote}
\emph{paṇḍuga rātri peddalu pēkāṭa āḍāru.}
\end{quote}

\begin{quote}
On the festival night the grown-ups played cards.
\end{quote}

\te{పేక}, playing cards, + \te{ఆట}, game.

\begin{cousinweb}
\textbf{\te{పేకాట}} (\emph{pēkāṭa}) is built from \te{పేక}, \textquotedblleft{}playing cards\textquotedblright{}, and native \te{ఆట}, \textquotedblleft{}play, game\textquotedblright{} (DEDR 347); the origin of \te{పేక} is not settled.
\end{cousinweb}

\begin{grammarlens}[title={What you've built}]
One more word for describing what you lived and felt.
\end{grammarlens}

\subsection*{Your turn}
\begin{itemize}
\raggedright
  \item \emph{Say it:} \emph{pēkāṭa}
  \item \emph{Say it:} \emph{pēkāṭa}, once more
  \item \emph{Say it:} \emph{paṇḍuga rātri peddalu pēkāṭa āḍāru.} — On the festival night the grown-ups played cards.
  \item \emph{Recall:} say the Telugu for gardening, then the Telugu for knitting, then say \emph{pēkāṭa} again
\end{itemize}

\begin{tcolorbox}[breakable,colback=teal!4,colframe=teal!35!black,title={Before you move on}]
What does \te{పేకాట} mean? (\textquotedblleft{}A card game\textquotedblright{}.) Say it once more.
\end{tcolorbox}
